% Fichier complété par tools/complete_missing_corrections.py.
% Source : DYN/DYN-06-PFD/08_RR3D/08_RR3D.tex
% Statut : rédigé (1 question(s)).
\begin{corrigebox}[Corrigé rédigé]
\CorrigeQuestion{1}
On écrit le théorème du moment dynamique sur l'axe indiqué, au
            point de concours de la liaison afin d'annuler les inconnues de réaction.
            Le membre droit est la projection de
            $\mathbf I\dot{\vec\Omega}+\vec\Omega\wedge(\mathbf I\vec\Omega)+
            \overrightarrow{AG}\wedge m\vec\Gamma_G$. Pour l'ensemble 1+2, les actions
            internes disparaissent. Les deux projections donnent le système différentiel
            couplé des angles $\theta$ et $\varphi$.
\end{corrigebox}
